% Document:    The Covering Number C(16,5,3) Is at Least 62
% Version:     v1.0.1
% Author:      Celaya Solutions
% Contact:     hello@celayasolutions.com
% Date:        2026-10-06
% SHA256:      11635a90ae72265bf483f9e6d13b4fc3bbff64a167832aad40939b6db32b9628
% Chain:       solana-mainnet
% Tx:          4SUk2RkFrqV8x9fJVBqpyTgRLyWD8r66Q4PRdUrBe3r5SymcGEh5HYzYGRByZGfuhHKzjqbR19uNVYks3Eyjq6uk
% License:     CC BY 4.0 / Celaya Solutions

\documentclass[11pt]{amsart}

\usepackage{amssymb}
\usepackage{booktabs}
\usepackage{array}
\usepackage{xurl}
\usepackage[hidelinks]{hyperref}
\hypersetup{pdftitle={The covering number C(16,5,3) is at least 62},
  pdfauthor={Christopher B. Celaya}}

\newtheorem{theorem}{Theorem}[section]
\newtheorem{lemma}[theorem]{Lemma}
\newtheorem{proposition}[theorem]{Proposition}
\newtheorem{corollary}[theorem]{Corollary}
\theoremstyle{definition}
\newtheorem{definition}[theorem]{Definition}
\theoremstyle{remark}
\newtheorem{remark}[theorem]{Remark}

\newcommand{\cB}{\mathcal{B}}
\newcommand{\cD}{\mathcal{D}}
\newcommand{\cR}{\mathcal{R}}

\begin{document}

\title{The covering number \texorpdfstring{$C(16,5,3)$}{C(16,5,3)} is at least 62}

\author{Christopher B. Celaya}
\address{Celaya Solutions Research, El Paso, Texas, USA}
\email{hello@celayasolutions.com}

\subjclass[2020]{Primary 05B40; Secondary 05B05, 68V05, 68V15}
\keywords{covering design, covering number, Sch\"onheim bound, Tur\'an number,
SAT solving, DRAT proof, computer-assisted proof}

\begin{abstract}
A $(v,k,t)$ covering is a family of $k$-subsets (blocks) of a $v$-set such
that every $t$-subset lies in at least one block, and the covering number
$C(v,k,t)$ is the least number of blocks in such a family. The Sch\"onheim
bound gives $C(16,5,3)\ge 61$, and a covering with 65 blocks has been known
since 1997; these were the best recorded bounds. We prove that
$C(16,5,3)\ge 62$. Counting shows that a 61-block covering would be very
rigid: one point lies in 20 blocks and every other point in 19, every triple is
covered once or twice, the pairs covered six times form a five-edge star plus a
perfect matching of the remaining ten points, and the derived design at a point
of degree 19 is a minimum $(15,4,2)$ covering. Minimum $(15,4,2)$ coverings
fall into exactly four isomorphism classes, as Allston, Buskens and Stanton
showed; we confirm this by two independent exhaustive computations. Fixing the
derived design at a suitable point to each class in turn leaves four Boolean
satisfiability instances, each with 144,553 variables and 512,424 clauses. All
four are unsatisfiable: the SAT solver Lingeling refutes them with DRAT
proofs of 13.0 to 15.6 million lines that the proof checker drat-trim
verifies, and the solver CaDiCaL and a separately written constraint
programming model agree. The encoder, the
formulas, the proofs and the checking logs are publicly available.
\end{abstract}

\maketitle

\section{Introduction}\label{sec:intro}

Let $v\ge k\ge t\ge 1$ be integers. A \emph{$(v,k,t)$ covering} is a family
$\cB$ of $k$-subsets, called \emph{blocks}, of a $v$-set $V$ such that every
$t$-subset of $V$ is contained in at least one block. The \emph{covering
number} $C(v,k,t)$ is the least size of a $(v,k,t)$ covering. We refer to the
survey of Gordon and Stinson~\cite{handbook} for background, and to the La
Jolla Covering Repository~\cite{ljcr}, now continued by the Covering
Repository~\cite{coveringrepo}, for tables of bounds. Taking complements,
$C(v,k,t)$ equals the Tur\'an number $T(v,v-t,v-k)$, the least number of
$(v-k)$-subsets of a $v$-set such that every $(v-t)$-subset contains one of
them~\cite{gordon1995,handbook}; for Tur\'an numbers in general see the
survey~\cite{sidorenko1995}.

The basic lower bound is due to Sch\"onheim~\cite{schonheim1964}:
\[
C(v,k,t)\ \ge\ L(v,k,t) := \left\lceil \frac{v}{k}\, L(v-1,k-1,t-1)\right\rceil,
\qquad L(v,k,0)=1 .
\]
For $(v,k,t)=(16,5,3)$ it gives
$L(16,5,3)=\lceil \tfrac{16}{5}\lceil \tfrac{15}{4}\lceil
\tfrac{14}{3}\rceil\rceil\rceil=\lceil \tfrac{16}{5}\cdot 19\rceil=61$.
A covering with 65 blocks was found by Rade Beli\'c in 1997~\cite{ljcr}, and Bluskov
obtained the same bound by a construction from the Steiner system
$S(3,5,17)$~\cite[Corollary~2.3.18]{bluskov1997}. The bounds
$61\le C(16,5,3)\le 65$ are those recorded by the covering
repositories~\cite{ljcr,coveringrepo}. General improvements of the
Sch\"onheim bound give nothing more here: the hypotheses of the Mills--Mullin
improvement~\cite{millsmullin1992} fail for these parameters, and the bounds
of Horsley and Singh~\cite{horsley2018} do not exceed 61. We know of no
earlier lower bound above 61 for $C(16,5,3)$ or, equivalently, for
$T(16,13,11)$. Our main result improves the lower bound.

\begin{theorem}\label{thm:main}
There is no $(16,5,3)$ covering with $61$ blocks. Hence
$62\le C(16,5,3)\le 65$.
\end{theorem}

The proof has a human part and a computer part. In
Section~\ref{sec:structure}, elementary counting shows that a 61-block
covering is extremely constrained: one point lies in 20 blocks and every
other point in 19; the pairs covered six times rather than five form a graph
consisting of a five-edge star at the point of degree 20 and a perfect
matching of the remaining ten points; every triple is covered once or twice,
with a prescribed number of doubly covered triples through each pair; and the
derived design at a point of degree 19 is a minimum $(15,4,2)$ covering. Minimum $(15,4,2)$ coverings have 19 blocks and fall into
four isomorphism classes (Section~\ref{sec:link}). This classification is due
to Allston, Buskens and Stanton~\cite{allston1988}; since our argument
depends on it, we confirm it by two independent exhaustive computations. In
Section~\ref{sec:reduction} we fix the derived design at a suitable point to
one of the four representatives. Each of the four resulting problems becomes a
propositional formula in conjunctive normal form (CNF), constructed so that
every 61-block covering yields a satisfying assignment of one of them. In
Section~\ref{sec:computation} we report that all four formulas are
unsatisfiable. The refutations were produced by the SAT solver
Lingeling~\cite{lingeling} as DRAT proofs and verified by the proof checker
drat-trim~\cite{wetzler2014}; the SAT solver CaDiCaL~\cite{cadical} and a
separately written CP-SAT model~\cite{ortools} give the same answer.

Computer-assisted proofs of this kind are common in combinatorics.
The value $C(10,5,4)=51$ was determined by Margot~\cite{margot2003}, with
branch-and-cut and isomorphism pruning, and independently by Applegate, Rains
and Sloane~\cite{applegate2003}, who also determined several further
covering numbers. SAT solvers with independently
checkable proofs settled the Boolean Pythagorean triples
problem~\cite{heule2016} and Keller's conjecture in dimension
seven~\cite{brakensiek2022}. Closest to the present paper,
Krug~\cite{krug2026} recently proved $C(12,6,4)=41$ by the same overall
strategy: counting forces the structure of a putative smaller covering, every
derived design is an optimal covering of smaller parameters, and the remaining
cases are refuted by a SAT solver with certified proofs.

All code, formulas, logs and proofs are public (Section~\ref{sec:data}).

\section{Notation and basic counting}\label{sec:prelim}

Throughout, $\cB$ is a family of 5-subsets of a 16-set $V$, a priori with
repetitions. For a point $x$, a pair $\{x,y\}$ and a triple $T$ we write
$r_x$, $\lambda_{xy}$ and $\mu_T$ for the numbers of blocks containing $x$,
$\{x,y\}$ and $T$, counted with multiplicity, and call $r_x$ the
\emph{degree} of $x$. The family
$\cB$ is a $(16,5,3)$ covering if and only if $\mu_T\ge 1$ for every triple
$T$. For a point $x$, the \emph{derived design} at $x$ is the family
(with multiplicities)
\[
\cD_x=\{B\setminus\{x\}: x\in B\in\cB\}
\]
of $r_x$ quadruples on $V\setminus\{x\}$. A pair $\{u,w\}\subseteq
V\setminus\{x\}$ lies in exactly $\mu_{xuw}$ members of $\cD_x$, and a point
$u\ne x$ lies in exactly $\lambda_{xu}$ of them. In particular, if $\cB$ is a
$(16,5,3)$ covering then $\cD_x$ is a $(15,4,2)$ covering. Two families are
\emph{isomorphic} if a bijection between their point sets carries one onto the
other.

\begin{lemma}\label{lem:basic}
Let $\cB$ be a $(16,5,3)$ covering. Then $\lambda_{xy}\ge 5$ for every pair
and $r_x\ge 19$ for every point.
\end{lemma}

\begin{proof}
Each of the 14 points $w\notin\{x,y\}$ needs a block containing
$\{x,y,w\}$, and a block through $\{x,y\}$ contains only three further points,
so $\lambda_{xy}\ge\lceil 14/3\rceil=5$. Each block through $x$ contains four
pairs through $x$, so $4r_x=\sum_{y\ne x}\lambda_{xy}\ge 15\cdot 5=75$ and
$r_x\ge 19$.
\end{proof}

Since $\sum_x r_x=5|\cB|$, Lemma~\ref{lem:basic} gives $5|\cB|\ge 16\cdot 19$,
that is $|\cB|\ge 61$; this is the Sch\"onheim bound. In particular a 61-block
covering has no repeated block, since deleting one copy of a repeated block
would leave a covering with 60 blocks.

\section{The structure of a 61-block covering}\label{sec:structure}

In this section $\cB$ is a $(16,5,3)$ covering with exactly 61 blocks; by the
remark above, its blocks are distinct.

\begin{lemma}\label{lem:degrees}
There is a point $z$ with $r_z=20$, and $r_x=19$ for every $x\ne z$.
\end{lemma}

\begin{proof}
The degrees sum to $5\cdot 61=305=16\cdot 19+1$ and each is at least 19 by
Lemma~\ref{lem:basic}.
\end{proof}

For a pair $\{x,y\}$ put $e_{xy}=\lambda_{xy}-5\ge 0$. Then
\begin{equation}\label{eq:pairexcess}
\sum_{y\ne x} e_{xy}=4r_x-75=
\begin{cases} 1 & \text{if } r_x=19,\\ 5 & \text{if } x=z.\end{cases}
\end{equation}

\begin{lemma}\label{lem:pairs}
Every pair has $\lambda_{xy}\in\{5,6\}$. Let $E$ be the graph on $V$ whose
edges are the pairs with $\lambda_{xy}=6$. Then there are a $5$-set
$A\subseteq V\setminus\{z\}$ and a perfect matching $P$ of
$M=V\setminus(A\cup\{z\})$ such that
\[
E=\bigl\{\{z,a\}:a\in A\bigr\}\cup P .
\]
In particular every point other than $z$ has exactly one neighbour in $E$, and
$E$ has ten edges and no triangle.
\end{lemma}

\begin{proof}
By \eqref{eq:pairexcess}, a point $x\neq z$ has $e_{xy}=1$ for exactly one $y$
and $e_{xy}=0$ for all other $y$. In particular $e_{zy}\le 1$ for every
$y\ne z$, so the five units of excess at $z$ lie on five pairs $\{z,a\}$ with
$a$ in a $5$-set $A$. Each $a\in A$ has used its single unit on $\{z,a\}$, so
it has no other neighbour in $E$. Every $x\in M$ therefore has its unique
neighbour in $M$, and these pairs form a perfect matching of $M$.
\end{proof}

For a pair $\{x,y\}$, each block through $\{x,y\}$ contains three triples
through $\{x,y\}$, so $\sum_{w}\mu_{xyw}=3\lambda_{xy}$, where $w$ runs over
the 14 other points. Hence
\begin{equation}\label{eq:tripleexcess}
\sum_{w\notin\{x,y\}}(\mu_{xyw}-1)=3\lambda_{xy}-14=
\begin{cases} 1 & \text{if } \{x,y\}\notin E,\\ 4 & \text{if }
\{x,y\}\in E.\end{cases}
\end{equation}
Call a triple \emph{doubled} if $\mu_T=2$.

\begin{lemma}\label{lem:triples}
Every triple $T$ has $\mu_T\in\{1,2\}$. Each pair outside $E$ lies in exactly
one doubled triple, each pair in $E$ lies in exactly four doubled triples, and
there are exactly $50$ doubled triples.
\end{lemma}

\begin{proof}
If $\mu_T\ge 3$, then $T$ contributes at least $2$ to the left side of
\eqref{eq:tripleexcess} for each of its three pairs. A pair outside $E$ has
total $1$ there, so all three pairs of $T$ lie in $E$; but $E$ has no
triangle. So $\mu_T\le 2$ for all $T$, and \eqref{eq:tripleexcess} counts the
doubled triples through each pair. Let $D$ be the number of doubled
triples. Counting incidences between pairs and doubled triples gives
$3D=110\cdot 1+10\cdot 4$, so $D=50$. (As a check,
$10\cdot 61=\binom{16}{3}+D$.)
\end{proof}

\begin{lemma}\label{lem:derived}
Let $x\ne z$ and let $p$ be the neighbour of $x$ in $E$. Then $\cD_x$ is a
$(15,4,2)$ covering with $19$ distinct blocks, in which $p$ lies in six blocks
and every other point in five. The pairs that lie in two blocks of $\cD_x$ are
the pairs $\{u,w\}$ for which $\{x,u,w\}$ is doubled; they form a star with
centre $p$ and four leaves together with a perfect matching of the remaining
ten points, and every other pair lies in exactly one block of $\cD_x$.
\end{lemma}

\begin{proof}
$\cD_x$ has $r_x=19$ blocks, and they are distinct because the blocks of $\cB$
are. The degree of $u$ in $\cD_x$ is $\lambda_{xu}$, which is $6$ for $u=p$
and $5$ otherwise. The multiplicity of $\{u,w\}$ in $\cD_x$ is $\mu_{xuw}\in
\{1,2\}$. By Lemma~\ref{lem:triples}, $\{x,p\}$ lies in four doubled triples
and every other pair $\{x,u\}$ in exactly one, so in the graph of pairs
covered twice by $\cD_x$ the point $p$ has degree four and every other point
degree one. Such a graph is a star at $p$ with four leaves together with a
perfect matching of the other ten points.
\end{proof}

We record what the reduction in Section~\ref{sec:reduction} uses.

\begin{proposition}\label{prop:necessary}
If $\cB$ is a $61$-block $(16,5,3)$ covering, then its blocks are distinct
and, with $z$, $A$, $M$ and $E$ as above:
\begin{enumerate}
\item[(N1)] every triple $T$ satisfies $1\le\mu_T\le 2$;
\item[(N2)] every pair outside $E$ lies in exactly one doubled triple, and
every pair in $E$ in exactly four;
\item[(N3)] $E$ consists of the five pairs $\{z,a\}$, $a\in A$, and a perfect
matching of $M$;
\item[(N4)] for every $a\in A$, the derived design $\cD_a$ is a $(15,4,2)$
covering with $19$ blocks in which $z$ lies in six blocks.
\end{enumerate}
\end{proposition}

\begin{remark}\label{rem:xz}
The counting gives more, although we do not use it. The doubled triples
through $z$ form a graph $X_z$ on $V\setminus\{z\}$ in which every point of
$A$ has degree $4$ and every point of $M$ has degree $1$. If $X_z$ has $n_2$
edges inside $A$, it has $20-2n_2$ edges between $A$ and $M$ and $n_2-5$ edges
inside $M$, so $5\le n_2\le 10$.
\end{remark}

\section{Minimum \texorpdfstring{$(15,4,2)$}{(15,4,2)} coverings}\label{sec:link}

It is classical that $C(15,4,2)=19$~\cite{mills1972,mills1973}; the
Sch\"onheim bound is $\lceil \tfrac{15}{4}\lceil\tfrac{14}{3}\rceil\rceil=19$.
The argument of Lemma~\ref{lem:derived} applies to every minimum
$(15,4,2)$ covering.

\begin{lemma}\label{lem:link-structure}
Let $\cD$ be a $(15,4,2)$ covering with $19$ blocks, a priori with
repetitions. Then one point $h$, the \emph{hub}, lies in six blocks and every
other point in five. No pair lies in three blocks. The pairs lying in two
blocks form a star with centre $h$ and four leaves together with a perfect
matching of the other ten points, and every other pair lies in exactly one
block. In particular the blocks of $\cD$ are distinct.
\end{lemma}

\begin{proof}
Each point lies in at least $\lceil 14/3\rceil=5$ blocks, and the degrees sum
to $76=15\cdot 5+1$, so there is exactly one point of degree $6$. For a point
$u$ of degree $d_u$ we have $\sum_{w\ne u}(m_{uw}-1)=3d_u-14$, where $m_{uw}$
is the number of blocks containing $\{u,w\}$; this is $4$ at the hub and $1$
elsewhere. A pair in three blocks would contribute $2$ at each of its ends,
which is impossible because at most one end is the hub. So all $m_{uw}\le 2$,
and in the graph of pairs with $m_{uw}=2$ the hub has degree four and every
other point degree one; such a graph is a star at the hub with four leaves
together with a perfect matching of the other ten points. A repeated block would put
all six of its pairs in two blocks, but the graph of such pairs has no
triangle.
\end{proof}

We call the graph of pairs lying in two blocks the \emph{excess graph} of
$\cD$; it is isomorphic to $K_{1,4}\cup 5K_2$. A 19-block $(15,4,2)$
covering on $\{2,\dots,16\}$ is in \emph{standard form} if its hub is $2$, the
leaves of its star are $3,4,5,6$, and its matching is
$\{7,8\},\{9,10\},\{11,12\},\{13,14\},\{15,16\}$. Let $G$ be the group of
permutations of $\{2,\dots,16\}$ that preserve this standard excess graph; it
has order $4!\cdot 5!\cdot 2^5=92{,}160$. Every 19-block
$(15,4,2)$ covering is isomorphic to one in standard form. An isomorphism
between two coverings preserves pair multiplicities and hence excess graphs,
so an isomorphism between two coverings in standard form is an element of $G$.
Isomorphism classes of 19-block $(15,4,2)$ coverings therefore correspond to
$G$-orbits of coverings in standard form.

\begin{theorem}\label{thm:classification}
There are exactly four isomorphism classes of $(15,4,2)$ coverings with $19$
blocks. Representatives $\cR_1,\dots,\cR_4$ in standard form are listed in
Table~\ref{tab:reps}; their automorphism groups have orders $4$, $2$, $2$
and $4$.
\end{theorem}

That there are exactly four classes is due to Allston, Buskens and
Stanton~\cite[Theorem~12]{allston1988}, who call these designs D2, D1, B4 and
B3 and report the same automorphism group orders. Explicit isomorphisms from
their published block lists onto $\cR_1,\dots,\cR_4$ are in the
repository~\cite{covering64repo}. Their classification used computer
isomorphism testing. Only the first assertion of
Theorem~\ref{thm:classification}, that every 19-block $(15,4,2)$ covering is
isomorphic to some $\cR_i$, is used in the proof of Theorem~\ref{thm:main}.
Since the proof rests on it, we proved Theorem~\ref{thm:classification} again
by two exhaustive computations, written separately and sharing no code; each
of Computations~1 and~2 below is a complete proof.

\subsection*{Computation 1: hub configurations}
Let $\cD$ be in standard form. The hub lies in six \emph{hub blocks}, which
have $18$ further slots. Each leaf fills two of them, in different hub blocks,
and each of the ten matched points fills one. Regard the six hub blocks as
vertices; each leaf gives an edge of a graph $\Gamma_L$ joining the two hub
blocks containing it, and each matching pair $\{m,m'\}$ gives an edge of a
multigraph $\Gamma_M$ joining the hub blocks containing $m$ and $m'$ (a loop if
they are equal). Every vertex has total degree three, a loop counting twice,
and $\Gamma_L$ is simple, since two leaves in the same two hub blocks would
cover their pair twice. Two hub-block sets lie in the same $G$-orbit exactly
when their coloured graphs $(\Gamma_L,\Gamma_M)$ are isomorphic. There are
$1335$ labelled graphs $\Gamma_L$ in eight isomorphism classes, giving $206$
classes of coloured graphs $(\Gamma_L,\Gamma_M)$. For a representative of each class,
the remaining $13$ blocks avoid the hub and must cover a prescribed multigraph
of residual pair demands exactly. A depth-first search that branches on the
pair with fewest remaining candidate blocks enumerated every such
decomposition; all $206$ searches completed, with $31{,}989$ nodes in total.
Exactly four classes extend, with $12$, $4$, $2$ and $96$ completions; the
other $202$ have none. A separate program replayed every search tree and
verified each branching, forced move and dead end, and a CP-SAT enumeration of
the four feasible cases returned the same $114$ completions. Within each
feasible class, with hub-block set $H$, any isomorphism between completions
lies in the stabilizer $G_H$ of $H$ in $G$; these stabilizers have orders $48$, $8$, $4$ and
$384$, and in each class the completions form a single $G_H$-orbit. Hence
there are exactly four isomorphism classes, with automorphism groups of orders
$48/12=4$, $8/4=2$, $4/2=2$ and $384/96=4$.

\subsection*{Computation 2: direct enumeration}
A second program, written from scratch without reference to the first,
enumerates all $19$-block coverings in standard form directly, as an exact
cover problem in which the nine excess pairs must be covered twice and the
other $96$ pairs once. Two depth-first searches with different branching rules
both found the same $N=138{,}240$ coverings. As a third check, the program
listed all $5{,}745{,}600$ labelled hub-block sets, which fall into $206$
$G$-orbits, and for one representative $H$ of each orbit CP-SAT enumerated
all completions of $H$; they were exactly the coverings with hub-block set $H$
found by the depth-first searches. Since $G$ maps the completions of $H$ onto
those of $gH$, this gives again $N=\sum_H |G\cdot H|\,c(H)=1920\cdot 12+
11520\cdot 4+23040\cdot 2+240\cdot 96=138{,}240$, where $c(H)$ is the number
of completions of $H$. Applying
all $92{,}160$ elements of $G$ partitions these coverings into four orbits, of
sizes $23{,}040$, $46{,}080$, $46{,}080$ and $23{,}040$, so the automorphism
groups have orders $4$, $2$, $2$ and $4$, and
$N=\sum_i |G|/|\mathrm{Aut}(\cR_i)|$. The four representatives in
Table~\ref{tab:reps} lie in the four different orbits.

\begin{table}[ht]
\centering
\small
\caption{Representatives of the four classes of 19-block $(15,4,2)$
coverings, in standard form (hub $2$, leaves $3,4,5,6$, matching
$\{7,8\},\dots,\{15,16\}$); each entry is a block. In the notation of
\cite{allston1988} these are D2, D1, B4 and B3, and in the accompanying
repository \texttt{shape-1}, \texttt{shape-4}, \texttt{shape-44} and
\texttt{shape-47}.}\label{tab:reps}
\begin{tabular}{@{}llll@{}}
\toprule
$\cR_1$ & $\cR_2$ & $\cR_3$ & $\cR_4$\\
\midrule
2 3 4 5   & 2 3 4 5   & 2 3 4 5   & 2 3 4 5\\
2 3 6 7   & 2 3 6 7   & 2 3 6 7   & 2 3 7 8\\
2 4 6 8   & 2 4 6 8   & 2 4 9 11  & 2 4 9 10\\
2 5 9 10  & 2 5 9 11  & 2 5 10 13 & 2 5 11 12\\
2 11 13 15& 2 10 13 15& 2 6 15 16 & 2 6 13 14\\
2 12 14 16& 2 12 14 16& 2 8 12 14 & 2 6 15 16\\
3 8 11 12 & 3 8 13 14 & 3 8 9 10  & 3 6 9 10\\
3 9 13 16 & 3 9 10 16 & 3 11 12 15& 3 11 13 14\\
3 10 14 15& 3 11 12 15& 3 13 14 16& 3 12 15 16\\
4 7 13 14 & 4 7 15 16 & 4 6 10 12 & 4 6 11 12\\
4 9 12 15 & 4 9 12 13 & 4 7 8 16  & 4 7 13 15\\
4 10 11 16& 4 10 11 14& 4 13 14 15& 4 8 14 16\\
5 6 15 16 & 5 6 10 12 & 5 6 9 14  & 5 6 7 8\\
5 7 11 12 & 5 7 13 14 & 5 7 8 15  & 5 9 13 16\\
5 8 13 14 & 5 8 15 16 & 5 11 12 16& 5 10 14 15\\
6 9 11 14 & 6 9 14 15 & 6 8 11 13 & 7 9 12 14\\
6 10 12 13& 6 11 13 16& 7 9 12 13 & 7 10 11 16\\
7 8 9 10  & 7 8 9 10  & 7 10 11 14& 8 9 11 15\\
7 8 15 16 & 7 8 11 12 & 9 10 15 16& 8 10 12 13\\
\midrule
$|\mathrm{Aut}|=4$ & $|\mathrm{Aut}|=2$ & $|\mathrm{Aut}|=2$ &
$|\mathrm{Aut}|=4$\\
\bottomrule
\end{tabular}
\end{table}

\section{Reduction to four satisfiability problems}\label{sec:reduction}

Suppose that $\cB$ is a 61-block $(16,5,3)$ covering, and let $z$, $A$, $M$,
$E$ be as in Section~\ref{sec:structure}. Choose $a\in A$. By
Proposition~\ref{prop:necessary}(N4), $\cD_a$ is a 19-block $(15,4,2)$
covering with hub $z$. By Theorem~\ref{thm:classification} there are an index
$i\in\{1,2,3,4\}$ and a bijection $\sigma\colon V\setminus\{a\}\to
\{2,\dots,16\}$ that carries $\cD_a$ onto $\cR_i$. It carries the hub $z$ to
the hub $2$ of $\cR_i$. Extend $\sigma$ by $\sigma(a)=1$ and replace $\cB$ by
its image. After this relabeling:
\begin{itemize}
\item $a=1$ and $z=2$, and the blocks through $1$ are exactly the 19 blocks
$\{1\}\cup R$ with $R\in\cR_i$;
\item $A=\{1\}\cup A'$ for a 4-subset $A'\subseteq\{3,\dots,16\}$, and $E$
consists of $\{1,2\}$, the pairs $\{2,v\}$ with $v\in A'$, and a perfect
matching of $\{3,\dots,16\}\setminus A'$;
\item conditions (N1) and (N2) hold.
\end{itemize}
What is unknown is $A'$, the matching, and the 42 blocks that avoid the
point $1$.

\subsection*{The formula \texorpdfstring{$F_i$}{F\_i}}
For each $i$ we define a CNF formula $F_i$ with the following variables:
$x_B$ for each 5-subset $B$ of $\{1,\dots,16\}$ ($4368$ variables, meaning
that $B$ is a block); $d_T$ for each 3-subset $T$ ($560$ variables, meaning
that $T$ is doubled); $a_v$ for $v\in\{3,\dots,16\}$ ($14$ variables, meaning
that $v\in A'$); and $m_{uv}$ for each pair $\{u,v\}\subseteq\{3,\dots,16\}$
($91$ variables, meaning that $\{u,v\}$ is a matching edge). For a pair $Q$
let $\varepsilon_Q$ be the constant or literal expressing $Q\in E$:
\[
\varepsilon_{\{1,2\}}=\top,\quad \varepsilon_{\{1,v\}}=\bot,\quad
\varepsilon_{\{2,v\}}=a_v,\quad \varepsilon_{\{u,v\}}=m_{uv}
\quad (3\le u,v\le 16).
\]
The constraints are:
\begin{enumerate}
\item[(C1)] $x_B$ is true for the 19 sets $B=\{1\}\cup R$ with $R\in\cR_i$,
and false for every other 5-set containing $1$;
\item[(C2)] for every triple $T$: $1\le\sum_{B\supseteq T}x_B\le 2$, and $d_T$
is true if and only if $\sum_{B\supseteq T}x_B\ge 2$;
\item[(C3)] exactly four of the variables $a_v$ are true;
\item[(C4)] $m_{uv}\to\neg a_u$ and $m_{uv}\to\neg a_v$; every
$v\in\{3,\dots,16\}$ has at most one $u$ with $m_{uv}$ true; and $a_v$ or
some $m_{uv}$ is true;
\item[(C5)] for every pair $Q$: $\sum_{T\supseteq Q}d_T\ge 1$; if
$\varepsilon_Q$ holds then $\sum_{T\supseteq Q}d_T=4$, and otherwise
$\sum_{T\supseteq Q}d_T=1$.
\end{enumerate}
The cardinality conditions in (C2) and (C5) are encoded with a two-sided
variant of the sequential counter of Sinz~\cite{sinz2005}, in which each
counter variable is equivalent to the statement that at least $j$ of the
first $\ell$ inputs are true: for each triple $T$ a counter with three
registers runs over the $78$ variables $x_B$, $B\supseteq T$, and for each
pair $Q$ a counter with five registers runs over the $14$ variables $d_T$,
$T\supseteq Q$. Condition (C3) uses the sequential-counter encoding of
PySAT~\cite{ignatiev2018}, whose auxiliary variables we do not interpret, and
the at-most-one conditions in (C4) use the pairwise encoding. In (C5), when
$\varepsilon_Q$ is a literal, the two cases are implications guarded by
$\varepsilon_Q$ and by $\neg\varepsilon_Q$. Each $F_i$ has $144{,}553$
variables and $512{,}424$ clauses (Table~\ref{tab:families}). The four
formulas differ only in the unit clauses of (C1).

\begin{table}[ht]
\centering
\small
\caption{Variables and clauses of each formula $F_i$ by constraint family.
All counter variables are auxiliary.}\label{tab:families}
\begin{tabular}{@{}lrr@{}}
\toprule
Family & Variables & Clauses\\
\midrule
Block variables $x_B$; (C1) as unit clauses & 4{,}368 & 1{,}365\\
Doubled-triple variables $d_T$ & 560 & 0\\
(C2) triple counters and their outputs & 131{,}040 & 478{,}800\\
Variables $a_v$; (C3) with 80 auxiliary variables & 94 & 160\\
(C4) matching variables and clauses & 91 & 1{,}288\\
(C5) pair counters and their outputs & 8{,}400 & 30{,}811\\
\midrule
Total & 144{,}553 & 512{,}424\\
\bottomrule
\end{tabular}
\end{table}

\begin{proposition}\label{prop:reduction}
If a $(16,5,3)$ covering with $61$ blocks exists, then at least one of the
formulas $F_1,\dots,F_4$ is satisfiable.
\end{proposition}

\begin{proof}
Relabel $\cB$ as above, with derived design $\cR_i$ at the point $1$. Set
$x_B$ true exactly for the blocks of $\cB$, $d_T$ true exactly for the doubled
triples, $a_v$ true exactly for $v\in A'$, $m_{uv}$ true exactly for the
matching edges, every counter variable of (C2) and (C5) to the truth value
of the statement it represents, and the auxiliary variables of (C3) to values
satisfying its clauses, which exist because exactly four $a_v$ are true
(checked exhaustively; Section~\ref{sec:computation}). Condition (C1) holds
because the blocks of $\cB$ are distinct and the blocks through $1$ are the
sets $\{1\}\cup R$, $R\in\cR_i$;
(C2) is (N1) together with the definition of a doubled triple; (C3) holds
because $|A'|=4$; (C4) holds because the points of $A'$ have no neighbour in
$E$ other than $2$, while every point of $\{3,\dots,16\}\setminus A'$ has
exactly one neighbour in $E$, and it lies in $\{3,\dots,16\}\setminus A'$.
For (C5), the description of $E$ after relabeling shows that $\varepsilon_Q$
is true exactly when $Q\in E$: the only neighbour of $1$ in $E$ is $2$, the
neighbours of $2$ are $1$ and the points of $A'$, and the edges of $E$ inside
$\{3,\dots,16\}$ are the matching edges. Then (C5) is (N2).
\end{proof}

\begin{remark}\label{rem:exact}
Neither the number of blocks nor the pair multiplicities are written into
$F_i$. They are implied, so the reduction loses nothing: in any satisfying
assignment, (C3) and (C4) give a graph $E$ with $1+4+5=10$ edges, so (C5)
gives $3D=110+40$ for the number $D$ of doubled triples, and (C2) gives
$10|\cB|=\sum_T\mu_T=560+D=610$, where $\cB$ is the set of $B$ with $x_B$
true. For a pair $Q$, (C2) and (C5) give $3\lambda_Q=\sum_{w\notin
Q}\mu_{Q\cup\{w\}}=14+1+3[\varepsilon_Q]$, so $\lambda_Q=5+[\varepsilon_Q]$;
hence $4r_2=15\cdot 5+5$ and $r_2=20$. Conversely, if a 61-block covering has
the sets $\{1\}\cup R$, $R\in\cR_i$, as its blocks through $1$ and its point
of degree 20 is $2$, then $\lambda_{12}=6$ because $2$ is the hub of $\cR_i$,
so $1\in A$, and the proof of Proposition~\ref{prop:reduction} applies with
$a=1$ and $\sigma$ the identity. Thus $F_i$ is satisfiable if and only if
there is a 61-block covering whose blocks through $1$ are the sets $\{1\}\cup
R$, $R\in\cR_i$, and whose point of degree 20 is $2$.
\end{remark}

\section{Computation and certification}\label{sec:computation}

\begin{theorem}\label{thm:unsat}
The formulas $F_1,F_2,F_3,F_4$ are unsatisfiable.
\end{theorem}

We recall the certification method. A DRAT proof of unsatisfiability of a
CNF formula is a list of clause additions and deletions ending with the empty
clause. Each added clause, or \emph{lemma}, must be a reverse unit
propagation (RUP) consequence of the clauses present at that point (assigning
all its literals false and applying unit propagation yields a conflict), or
more generally a resolution asymmetric tautology (RAT). Both kinds of addition
preserve satisfiability and clause deletions also preserve satisfiability,
so a valid proof that reaches the empty clause shows that the
formula is unsatisfiable. In backward mode, drat-trim starts from the empty
clause and checks only the lemmas needed to derive it; these lemmas and the
original clauses they use form the \emph{core}. DIMACS is the standard text
format for CNF formulas.

The formulas were generated by a short Python script using PySAT
(version 1.9.dev15)~\cite{ignatiev2018} and written in DIMACS format. Each was
solved by Lingeling (version \texttt{bbc-9230380-\allowbreak 160707}, the SAT
Competition 2016 version~\cite{lingeling}, as bundled with PySAT), which
recorded a proof of unsatisfiability in the DRAT format. Each proof was then
checked by drat-trim~\cite{wetzler2014} (commit \texttt{2e3b2dc}) in backward
mode with the time limit raised (option \texttt{-t 1000000}), which reported
\texttt{s VERIFIED} in all four cases. Table~\ref{tab:runs} gives the
figures. drat-trim reports no RAT lemmas in any core, so every lemma in the
verified cores is a RUP lemma. The runs were
made on 5~October~2026 on one Apple M5 Max laptop with 18 cores and 128\,GB
of memory; the four cases ran concurrently with other jobs, so the times are
wall-clock times on a shared machine and are given only for orientation.

\begin{table}[ht]
\centering
\small
\caption{Refutation of $F_1,\dots,F_4$ by Lingeling and verification of the
DRAT proofs by drat-trim. Times in seconds; sizes in decimal gigabytes.
Solve is the wall-clock time of the Lingeling call through PySAT, including
loading the clauses and retrieving the proof; check times and core figures
are as reported by drat-trim.}\label{tab:runs}
\begin{tabular}{@{}lrrrrrr@{}}
\toprule
 & Solve & Proof lines & Size & Check & Core clauses & Core lemmas\\
\midrule
$F_1$ & 14{,}714 & 13{,}065{,}165 & 7.04 & 5{,}269 & 223{,}884 & 1{,}914{,}196\\
$F_2$ & 19{,}495 & 15{,}609{,}452 & 8.82 & 3{,}512 & 228{,}423 & 2{,}143{,}298\\
$F_3$ & 19{,}640 & 15{,}471{,}588 & 8.48 & 3{,}439 & 235{,}146 & 2{,}179{,}394\\
$F_4$ & 14{,}774 & 13{,}044{,}849 & 7.30 & 4{,}829 & 230{,}704 & 1{,}765{,}503\\
\bottomrule
\end{tabular}
\end{table}

\subsection*{Cross-checks}
Two further computations agree with Theorem~\ref{thm:unsat}; neither is part
of the proof. First, CaDiCaL~1.9.5~\cite{cadical}, also called through PySAT
on the same formulas, reported unsatisfiability in $1{,}045$ to $1{,}255$
seconds after $3.5$ to $4.2$ million conflicts; these runs produced no
proofs. Second, a separately written CP-SAT model in Google OR-Tools
9.15~\cite{ortools} states the same conditions directly as linear constraints
($\sum_{B\supseteq T}x_B=1+d_T$ for every triple, $\sum_{T\supseteq Q}d_T=
1+3\varepsilon_Q$ and, redundantly, $\sum_{B\supseteq Q}x_B=5+\varepsilon_Q$
for every pair, and the conditions on $A'$ and the matching). It reported all
four problems infeasible in $1{,}623$ to $2{,}089$ seconds. This model shares
the reduction and the representatives with the CNF formulas, but none of the
encoding: no counters, no cardinality encodings and no DIMACS output.

\subsection*{Checks on the formulas}
A proof checker certifies that a formula is unsatisfiable, not that the
formula says what it is meant to say. We therefore checked the encoding
separately. The DIMACS files are regenerated byte for byte from the
representatives alone by the script \texttt{gen\_cnf.py} in the ancillary
files, which consists of the encoding lines of the solve script without the
solver call; their SHA256 checksums are listed with the data. That each
counter variable is equivalent to its intended statement follows by induction
on the number of inputs; we also tested the counter code exhaustively for up
to $12$ inputs with bounds up to $6$, and for $14$ inputs with the bounds $3$,
$5$ and $6$ (the bound $5$ on $14$ inputs is the case used in (C5)): in every
case each assignment of the inputs extends uniquely, with every counter
variable equal to its intended value. The clauses of (C3) and (C4) were checked
exhaustively over all assignments of the $a_v$: the clauses of (C3), with
their auxiliary variables, are satisfiable exactly for the $1001$ assignments
with four $a_v$ true, and the solutions of (C4) on the $m_{uv}$ are exactly
the $945$ perfect matchings of the complement of $A'$ for each of the $1001$
choices of $A'$. Finally, for $1200$ random assignments of the main
variables, extended to the auxiliary variables (the counters canonically, the
auxiliary variables of (C3) by a satisfying extension when one exists), the violated
constraint instances of each family were compared with the violated instances
of (C1)--(C5); there were no mismatches. We also confirmed that every
representative $\cR_i$ is a 19-block $(15,4,2)$ covering in standard form with
hub $2$, as the definition of $\varepsilon_Q$ requires. The scripts for these
checks, with their output, are in the directory
\nolinkurl{experiments/2026-10-06/lb61-encoder-audit} of the
repository~\cite{covering64repo}.

\subsection*{What the proof relies on}
Theorem~\ref{thm:main} rests on four things. First, the arguments proved in
the text: Lemma~\ref{lem:basic}, the counting of Section~\ref{sec:structure},
Lemma~\ref{lem:link-structure} and the reduction of
Proposition~\ref{prop:reduction}. Second, the completeness part of
Theorem~\ref{thm:classification}. It has been obtained three times, by
Allston, Buskens and Stanton and by the two computations of
Section~\ref{sec:link}, which agree; the search trees of Computation~1 were
replayed by a separate program, but unlike the DRAT proofs they are not
checked by an established proof checker. Third, the correctness of the short
encoder that produces $F_1,\dots,F_4$, including the PySAT cardinality
encoding it calls for (C3), checked as described above. Fourth, the
correctness of drat-trim. The proof does not rely on the correctness of
Lingeling, CaDiCaL or CP-SAT. The DRAT proofs could also be converted to the
LRAT format, in which each lemma lists the clauses that justify it, and
checked with a formally verified checker such as
cake\_lpr~\cite{tan2021cakelpr}, as was done in~\cite{krug2026}; we have not
yet done this.

\begin{proof}[Proof of Theorem~\ref{thm:main}]
By Proposition~\ref{prop:reduction}, a 61-block covering would make one of
$F_1,\dots,F_4$ satisfiable, contradicting Theorem~\ref{thm:unsat}. Hence
$C(16,5,3)\ge 62$. The covering of Beli\'c~\cite{ljcr} gives
$C(16,5,3)\le 65$; see also~\cite[Corollary~2.3.18]{bluskov1997}.
\end{proof}

\section{Concluding remarks}\label{sec:remarks}

The value of $C(16,5,3)$ is now known to be $62$, $63$, $64$ or $65$. Part of
the method carries over to more blocks. For any $(16,5,3)$ covering with $b$
blocks, the identities of Section~\ref{sec:structure} give
$\sum_{\{x,y\}}(\lambda_{xy}-5)=10b-600$ and $\sum_x(r_x-19)=5b-304$. At
$b=62$ at most six points have degree greater than $19$, so at least ten
points have degree $19$. The derived design at each of them is a $(15,4,2)$
covering with $19$ blocks, so by Theorem~\ref{thm:classification} it is
isomorphic to one of $\cR_1,\dots,\cR_4$. The difference is that the graph of
pairs with $\lambda_{xy}>5$ is no longer forced: at $62$ blocks the counting
allows many degree sequences and excess graphs, so an analogous case split
would be considerably larger.

\section{Data availability}\label{sec:data}

The repository~\cite{covering64repo}, licensed CC BY 4.0, contains the
structural argument; Computation~1 of Section~\ref{sec:link}
(\nolinkurl{experiments/2026-10-03/link-hub-independent}) and Computation~2
(\nolinkurl{experiments/2026-10-06/reclassify-15-4-2}) with their evidence; the
isomorphisms from the designs of~\cite{allston1988} onto
$\cR_1,\dots,\cR_4$ (\nolinkurl{experiments/2026-10-03/literature-and-novelty});
the encoder (\nolinkurl{experiments/2026-10-05/lb61-certified}) and its audit
(\nolinkurl{experiments/2026-10-06/lb61-encoder-audit}); the CaDiCaL and CP-SAT
cross-checks; and the solver records and drat-trim logs. The SAT
computations of Section~\ref{sec:computation} are those of commit
\texttt{fd7adc2}, and the classification and encoder checks added for this
paper are in commit \texttt{d490808}. The four CNF formulas and DRAT
proofs (about $32$\,GB uncompressed), with SHA256 checksums, the search trees
of Computation~1 and instructions for checking the proofs, are archived on
Zenodo~\cite{zenodo}. The ancillary files of this arXiv submission contain
the representatives $\cR_1,\dots,\cR_4$ (as the 19 blocks through the point
$1$), the script \texttt{gen\_cnf.py} that writes the four formulas without
solving them, the CP-SAT model and the SHA256 checksums of all formulas and
proofs.

\section*{Use of AI assistance}

AI coding assistants (OpenAI Codex and Anthropic Claude Code), working under
the author's direction, wrote the programs in the accompanying repository,
including both classification programs of Section~\ref{sec:link}, the
encoder and the CP-SAT model. They also helped develop the counting arguments
of Section~\ref{sec:structure}, search the literature, draft this manuscript
and check its references and numbers. In this paper, computations described
as independent or separately written are separate programs that share no
code; they were not written by different people. The author has checked the
arguments and the references and takes full responsibility for the content.

\bibliographystyle{amsplain}
\bibliography{refs}

\end{document}
